\documentclass[presentation,10pt]{beamer}
\usetheme{Madrid}
\usecolortheme{whale}

% --- Packages ---
\usepackage{amsmath}
\usepackage{amsfonts}
\usepackage{amssymb}
\usepackage{booktabs}
\usepackage{graphicx}

% --- Presentation Metadata ---
\title[Subsonic Marshak Wave Fitted Solutions]{Generalized Self-Similar Solutions of \\ Subsonic Marshak Waves}
\subtitle{Theory and Numerical Validation Structure}
\author[Ramot \& Krief]{Tal Ramot \and Shay I. Heizler \and Menahem Krief}
\institute[HUJI]{
  The Hebrew University of Jerusalem, Physics dept. 
}
\date{ICTT-29 \big| Aix-en-Provence \big| July 17, 2026}

\begin{document}

% --- Title Slide ---
\begin{frame}
  \titlepage
\end{frame}

% --- Slide 1: Motivation ---
\begin{frame}{Motivation}
  \begin{block}{Inertial Confinement Fusion (ICF)}
    \begin{itemize}
      \item Hohlraum Modeling Crucial for radiation temperature prediction and experimental measurement accuracy.
      \item Energy Delivery Key to analyzing laser--walls coupling and refining pulse shaping techniques.
      \item Historical Context Builds upon classical foundations established by \emph{Lindl (1995)}, \emph{Kauffman (1995)}, and \emph{Rosen (2005)}.
    \end{itemize}
  \end{block}

  \begin{alertblock}{Adding Shay \& Menahem's suggestions}
    \begin{itemize}
      \item Slides about the Aluminum article...
    \end{itemize}
  \end{alertblock}
\end{frame}

% --- Slide 2: General Introduction ---
\begin{frame}{The Theory Of Subsonic Heat Waves}
  \begin{alertblock}{Physics Of The Problem}
    \begin{itemize}
      \item We treat a semi-infinite wall subjected to high-temperature boundary heating.
      \item Dynamic equilibrium is mediated strictly via \textbf{photon transport}.
      \item In the \textbf{optically-thick limit}, an exceptionally sharp temperature front develops.
    \end{itemize}
  \end{alertblock}

  \vspace{0.3cm}
  \begin{block}{The Transfer Equation Coupled to Euler Equations}
    \begin{itemize}
      \item Coupled Boltzmann \& internal energy equation.
      \item The diffusion equation under the gray approximation.
      \item Euler equations for the hydrodynamic evolution of the material - in two words.
      \item The EOS for different materials fitted by H \& R (2003 or previously).
    \end{itemize}
  \end{block}
\end{frame}

\end{document}